\documentclass[]{article}

\usepackage{amsmath}
\usepackage{multirow}
\usepackage{natbib}
\usepackage{graphicx} 


\begin{document}

In this paper, we addressed the challenge of identifying weak gravitational lensing maps generated by anomalous hydrodynamical simulation physics, framing it as an Out-of-Distribution (OoD) detection problem. The core difficulty lies in distinguishing subtle differences in non-Gaussian structure from variations caused by known cosmological and baryonic nuisance parameters. Our proposed solution is a two-stage, data-driven framework designed to isolate these structural anomalies.

Our methodology first employs the Wavelet Scattering Transform (WST) to compress each noisy lensing map into a 217-dimensional feature vector, effectively summarizing its multi-scale morphological information. We then model the probability distribution of these features using a Conditional Normalizing Flow (CNF), conditioned on five physical parameters: $\Omega_m$, $S_8$, and three baryonic nuisance parameters. During inference, a Multi-Layer Perceptron regressor first estimates these parameters from a map's WST features. The final OoD score is then calculated as the negative log-likelihood of the features under the CNF, conditioned on the estimated parameters. This approach allows the model to marginalize over known physical variations and focus on identifying improbable structural features.

The results demonstrate the effectiveness of our framework. The parameter regression network successfully inferred the physical parameters from the WST features, with cosmological parameters being more constrained than baryonic ones, as expected. On a synthetic validation set designed to test the detection of structural deviations, our method achieved a partial Area Under the Curve of 0.2223 in the stringent false positive rate regime of 0.001 to 0.05, indicating a strong ability to identify anomalies with high confidence. When applied to the blind test set, the resulting distribution of OoD scores revealed a distinct, heavy-tailed population of maps with exceptionally high negative log-likelihoods, successfully flagging them as strong OoD candidates.

We have learned that the combination of WST feature extraction and conditional density estimation provides a powerful tool for tackling model uncertainty in weak lensing analyses. The WST proves capable of capturing the essential non-Gaussian statistics indicative of baryonic physics, while the conditional nature of the normalizing flow is crucial for disentangling true structural anomalies from expected physical variations. This work presents a principled method for identifying discrepancies between simulation codes, a critical step towards validating the models necessary for achieving the full scientific potential of upcoming cosmological surveys.

\end{document}
                